\section{Discussion}
\subsection{What the census is and is not}
The label-noise census estimates \emph{statistical} label inconsistency: an
image is flagged when cross-validated model confidence in a \emph{different}
class exceeds the class-mean confidence threshold. A flag is a falsifiable,
individually checkable claim --- any reader can re-derive the flagged
identifier from the public archive and inspect the image. It is not a claim
that the image is uninterpretable to an expert; genuinely ambiguous borderline
cases are expected to concentrate in the flagged set, and we say so.
\subsection{The admissibility gate}
Confident learning presumes predicted probabilities that separate the classes
reasonably well. When the out-of-fold model is weak (our synthetic check:
$<$10 gradient steps per sample), the estimator over-flags massively --- in an
early run on \texttt{breastmnist} it reported a 49.8\% noise rate that we
traced to a model with $\approx$9 total optimiser steps, i.e.\ an artefact,
not a property of the data. We therefore gate every census: out-of-fold
accuracy $\geq \max(0.60, \text{majority rate} + 0.05)$, and a minimum of
$\approx$250 optimiser steps per fold, scaling epochs to the dataset size.
Datasets failing the gate are reported as \emph{census inadmissible}; we
publish the gate failure rather than a garbage number. We regard this gate as
a necessary methodological contribution for anyone auditing small medical
benchmarks, where undertrained auditors are the norm, not the exception.
\subsection{Engineering for a two-core sandbox}
All training ran on 2 CPU cores with $\leq$2\,GB RAM. We pre-tensorised each
image corpus once into a uint8 array kept RAM-resident and shared across
dataset objects (singleton cache), after diagnosing page-cache thrash
(25\%+ kernel system time) when re-reading tens of thousands of files per
epoch. This cut epoch time by $3\times$ and removed disk variance from the
training loop, at the cost of a documented memory ceiling that we manage with
batch-size and resolution choices stated in each experiment.
\subsection{Uncertainty quantification}
Every accuracy in this paper carries a Wilson 95\,\% interval, computed
with \texttt{statsmodels} from the committed confusion matrices
(\texttt{results/confidence\_intervals.json}). Malaria: CNN
0.9608 [0.9529, 0.9675], RegionGCN 0.9598 [0.9518, 0.9665] ($n{=}2{,}758$);
the cleaned-retrain deltas sit well inside these intervals, so we report
the census-cleaning effect as \emph{small and statistically
indistinguishable from zero on this test set}, not as an improvement or a
regression. Pneumonia: CNN 0.7212 [0.6847, 0.7549], RegionGCN 0.7901
[0.7564, 0.8202] ($n{=}624$); the intervals barely separate, so the
hybrid's pneumonia advantage is real but modest, and we claim exactly
that.
\subsection{Interpretability check}
Integrated Gradients (Captum) on the committed malaria CNN shows the
model keys on parasite morphology: attribution hotspots coincide with
the visible parasite bodies on Parasitized cells, while Uninfected
cells draw diffuse, low-contrast attribution (Fig.~saliency). Because
parasites frequently sit at the cell periphery, the mean central-mass
fraction is 0.219 --- slightly \emph{below} the 0.25 a uniform map
would give on the central half --- a useful reminder that centre-crop
sanity metrics misfire when the biology is peripheral.
\subsection{Verification and robustness side-analyses}
Four independent checks (\texttt{results/tool\_battery\_1.json}):
(i) \textbf{sympy}: the GCN normalisation $D^{-1/2}(A{+}I)D^{-1/2}$ from
Methods is verified symbolically --- symmetric, spectrum in $({-1},1]$
($[-0.167,1]$ for $g{=}3$), and matching our NumPy implementation to
$10^{-10}$. (ii) \textbf{networkx}: the $3{\times}3$ region grid has
diameter 2 (two GCN layers mix every region pair), but the $4{\times}4$
pneumonia grid has diameter 3, so two layers leave the farthest region
pairs unmixed --- a stated architectural limitation, not a hidden one.
(iii) \textbf{OpenCV}: downsizing native-resolution cells to 48\,px with
INTER\_AREA vs INTER\_LINEAR moves the committed model's confidence by
0.025 on average and up to 0.253 on individual cells --- preprocessing
choice is a real deployment variable, so dxtool pins its resizer.
(iv) \textbf{seaborn}: OOF confidence distributions (Fig.~above) show
the one-directional noise tails directly.
\subsection{Source verification of the published baselines}
Benchmark numbers quoted from memory are a known failure mode, so both
published frontiers were re-derived from source
(\texttt{results/tool\_battery\_2.json}, \texttt{tool\_battery\_4.json}).
\textbf{NCBI E-utilities} and \textbf{Europe PMC} full-text search find
\emph{zero} prior label-noise censuses of either dataset (the novelty
claim). \textbf{CrossRef} verifies both benchmark DOIs. The Europe PMC
full text of Rajaraman et al.\ (PMC5907772), parsed with \textbf{lxml}
and independently re-parsed with \textbf{xmltodict} (identical numbers),
corrected two misquotes we initially carried: their custom CNN reports
$94.0 \pm 1.0$\% cell-level accuracy (the widely-quoted 95.9\% is the
\emph{patient}-level number, their Table~6), and the pre-trained
reference is VGG-16 at $94.5 \pm 1.5$\% --- no VGG-19 appears in the
paper. Our 96.08\% cell-level result therefore stands $+2.1$ points
above their cell-level custom CNN, not the $+0.2$ a memory-level quote
would suggest. Kermany's headline 92.8\%/93.2\%/90.1\% is paywalled;
all three numbers were corroborated against two independent open-access
citing papers (PMC10660141, PMC11759907) and are labelled as
secondary-sourced. \textbf{trafilatura} captured clean provenance text
for the NIH dataset record (PubMed/PMC/PeerJ/Mendeley landing pages are
browser-gated to automated fetches --- recorded as a limitation), and
\textbf{htmldate} dates the record to 2020-03-16. \textbf{BeautifulSoup}
extracted the LHC dataset table structure, confirming our release is the
thin-smear series.
\subsection{Pipeline integrity and consistency audits}
(\texttt{results/tool\_battery\_3.json}, \texttt{tool\_battery\_4.json}.)
\textbf{imageio} re-decoded 100 source files (60 malaria PNGs, 40
pneumonia JPEGs) through an independent code path and reproduced the
stored training arrays with \emph{zero} pixel difference --- the
preprocessing pipeline preserves the source bytes. \textbf{DuckDB} SQL
cross-tabs show pneumonia flags concentrate in the pneumonia class
(132 of 155), consistent with ambiguous infiltrate films being the noisy
ones. An adjusted \textbf{formulaic} logistic model of flag status on
the image properties shows pneumonia flags are driven by sharpness alone
(OR 3.09 per SD, $p{=}1.1{\times}10^{-17}$) while malaria flags are
driven by contrast (OR 1.58, $p{=}1.4{\times}10^{-4}$) and entropy
(OR 1.27, $p{=}0.013$) but \emph{not} sharpness --- and the low
pseudo-$R^2$ (0.017 malaria / 0.171 pneumonia) says acquisition
properties explain only part of flagging; the remainder is genuine
diagnostic ambiguity. \textbf{pdfplumber} re-extracted the rendered
paper's tables and confirmed every audited number matches the committed
JSONs; \textbf{PyMuPDF} renders pages for visual inspection.
\textbf{Plotly} and \textbf{openpyxl} produce the interactive
property-scatter and the flagged-image registry workbook shipped with
the results; \textbf{ReportLab} generates the one-page census review
card.

\subsection{Suite-level audit of the item 10a tabular panels}
The shared repository's top-level results cover 22 tabular diagnosis
datasets (the item 10a scope; 264 model--seed runs). We re-analysed them
with three independent toolchains
(\texttt{results/tena\_scope/tool\_battery\_6\_tena.json}).
\textbf{DuckDB} ranks the panels by cross-model difficulty: fertility
(best model AUC $\approx 0.68$) and haberman ($\approx 0.70$) remain
hard for every model tried; heart\_va, wpbc and ilpd follow.
\textbf{agate} recomputes the per-model aggregates and \textbf{csvkit}
(csvstat) profiles the flattened metrics a third time (panel ROC-AUC
mean 0.8429, median 0.8828) --- all three agree. A \textbf{statsmodels}
Benjamini--Hochberg FDR pass over per-dataset one-sided Wilcoxon tests
finds 0/22 significant, for a structural reason: with three seeds per
dataset the minimum achievable one-sided Wilcoxon $p$ is 0.25, so no
panel can reach FDR significance by design. Per-dataset superiority
claims for the 10a panels are therefore stated descriptively, not
inferentially. \textbf{leather} renders the per-dataset best-model AUC
chart (Fig.\ shipped as SVG).
\subsection{Flagged images are measurably atypical}
If the census flags a property of the images rather than an artefact of
the estimator, flagged images should differ measurably from unflagged
ones. Using scikit-image (variance-of-Laplacian sharpness, RMS contrast,
Shannon entropy) with Mann--Whitney tests against 500 matched controls
(\texttt{results/*/flagged\_image\_properties.json}): flagged pneumonia
films are significantly \emph{sharper} (median 0.0053 vs 0.0032,
$p{=}1.7{\times}10^{-27}$), higher-contrast ($p{=}4.9{\times}10^{-7}$)
and higher-entropy ($p{=}0.0035$) than controls --- acquisition-atypical
films are exactly where label errors concentrate. Flagged malaria cells
differ only in contrast ($p{=}0.0024$; sharpness $p{=}0.46$, entropy
$p{=}0.62$), consistent with subtler borderline parasitisation rather
than acquisition artefacts. The two diseases fail differently, and the
census sees both.
\subsection{Malaria census findings}
The census of the NIH malaria training partition (22{,}046 cells) estimates
a label-noise rate of 11.44\,\% (2{,}522 estimated errors; Table
\ref{tab:census}), with a strong directional asymmetry: 2{,}358 of the
estimated errors are cells labelled \emph{Parasitized} but cross-validated
as \emph{Uninfected}, against only 164 in the opposite direction. The
labels therefore err almost entirely in one direction, consistent with
over-inclusive parasitised annotations during dataset construction rather
than random annotator noise. We conservatively flag 249 high-confidence
cases (Appendix A). An independent cross-check with the published
\texttt{cleanlab} reference implementation flags 296 cases and confirms
\emph{all 249} of ours (Jaccard 0.841): our from-scratch estimator's
conservative flags are a strict subset of the reference library's, which
strengthens the claim that the flagged set is a property of the data, not
of our estimator.
\subsection{Pneumonia census findings}
The census of the Kermany training partition (4{,}708 films) estimates a
label-noise rate of 21.15\,\% (996 estimated errors), admissible under
the gate (OOF accuracy 0.898 against a 0.742 majority rate), with the
same one-directional signature as malaria but stronger: 973 images
labelled \emph{pneumonia} cross-validate as \emph{normal}, against only
23 in the opposite direction. We conservatively flag 155 high-confidence
cases (Appendix A), every one of which is independently flagged by the
\texttt{cleanlab} reference implementation (which flags 209; Jaccard
0.742). That two independently assembled public medical datasets exhibit
the same directional failure mode --- over-inclusive positive labels ---
is the central empirical finding of this study, and each flagged
identifier is published so any reader can inspect the films.
\subsection{Pneumonia: an honest gap, and where it lives}
Our from-scratch models underperform the published Kermany et al.\ number
(92.8\,\%, Inception-v3 transfer from ImageNet) by 13--21 accuracy points
(CNN 72.1\,\%, RegionGCN 79.0\,\%; Table \ref{tab:benchmark}). The gap
decomposes cleanly: sensitivity is 99\,\%+ on the official test set for
both models, while specificity is 45--55\,\% --- the models catch
pneumonia but over-call normals. Two causes are quantified in this study:
(i) the training split is 74\,\% pneumonia, so an unweighted loss rewards
over-calling (addressed in the tuned iteration by class-weighted loss);
(ii) the label-noise census (Section census results) measures how much of
the residual is annotation noise. We also note the comparison is
asymmetric by construction: Kermany et al.\ initialised from ImageNet
features, we train from random initialisation on 2 CPU cores; the honest
statement is that a from-scratch small model reaches 79\,\% on this
benchmark, and the published 92.8\,\% is not a from-scratch number.
\subsection{The hybrid's value is dataset-dependent}
The RegionGCN hybrid loses slightly to the CNN core on malaria
($-$0.10 pts) but wins clearly on pneumonia (+7.9 pts accuracy,
+1.8 pts AUC). Chest X-rays have large-scale spatial structure (heart
border, costophrenic angles) that region-level graph reasoning can
exploit; 48\,px cell crops do not. A hybrid head is therefore not a
uniform upgrade; it should be adopted per-dataset, with the delta
measured, exactly as reported here.
\subsection{Honest negatives}
\begin{itemize}
\item The RegionGCN hybrid did \emph{not} beat the plain CNN core on
malaria (95.98\,\% vs 96.08\,\% test accuracy; committed
\texttt{baseline\_results.json}). The graph head adds 18{,}464 parameters
and measurable training time for no gain on this task; we report it as a
negative result rather than tuning until it wins.
\item Our first breastmnist census produced a 49.8\,\% noise rate that was
an artefact of an undertrained auditor (see admissibility gate); the run
is preserved in the logs and motivated the gate.
\item RAM limits (2\,GB, no swap) forced malaria images to 48\,px; two
long training runs were killed by the kernel OOM killer before we adopted
per-fold checkpointing. Both mitigations are engineering compromises, not
scientific choices.
\end{itemize}

